% Part: normal-modal-logic
% Chapter: filtrations
% Section: euclidean-filtrations

\documentclass[../../../include/open-logic-section]{subfiles}

\begin{document}

\olfileid[pt]{nml}{fil}{euc}

\olsection{Filtrações de modelos euclidianos}

A abordagem de \olref[acc]{sec} não funciona para modelos euclidianos,
ou seriais e euclidianos. Considere o modelo na parte superior de
\olref{fig:ser-eucl}, que é tanto euclidiano quanto serial. Seja
$\Gamma = \{p, \Box p \}$. Ao tomar uma filtração por meio de $\Gamma$,
temos $[w_1] = [w_3]$, pois $w_1$ e $w_3$ são os únicos mundos que
concordam em~$\Gamma$. Toda filtração também terá a seta herdada de
$\mModel{M}$, como mostra \olref{fig:ser-eucl2}. Esse modelo não é
euclidiano. Além disso, não podemos acrescentar setas ao modelo para
torná-lo euclidiano. Teríamos de acrescentar setas duplas entre $[w_2]$
e $[w_4]$ e, depois, também entre $w_2$ e~$w_5$. Mas $\Box p$ deve ser
verdadeira em~$w_2$, ao passo que $p$ é falsa em~$w_5$.

\begin{figure}[htpb]
  \centering
  \begin{tikzpicture}[modal]
    \node[world] (w1)
          [label={left: \mFalse{p}},
            label={below: $\mSat{{}}{\Box p}$}] {$w_1$}; 
    \node[world] (w2)
          [label={right: \mTrue{p}},
            label = {below: $\mSat{{}}{\Box p}$},
            right=of w1] {$w_2$}; 
    \draw[->] (w1) to (w2);
    \draw[reflexive above] (w2) to (w2);  
    \node[world] (w3)
          [label={left: \mFalse{p}},
            label={below: $\mSat{{}}{\Box p}$},
            below=of w1] {$w_3$}; 
    \node[world] (w4)
          [label={right:\mTrue{p}},
            label={below: $\mSat/{{}}{\Box p}$},
            right=of w3] {$w_4$};
    \draw[->] (w3) to (w4);
    \draw[reflexive above] (w4) to (w4);     
    \node[world] (w5)
          [label={right: \mFalse{p}},
            label=below:$\mSat/{{}}{\Box p}$,
            right=of w4] {$w_5$}; 
    \draw[->, bend left] (w4) to (w5);
    \draw[->, bend left] (w5) to (w4);
    \draw[reflexive above] (w5) to (w5);     
  \end{tikzpicture}
  \caption{Um modelo serial e euclidiano.}
  \ollabel{fig:ser-eucl}
\end{figure}

\begin{figure}[ht]
  \centering
  \begin{tikzpicture}[modal,every text node part/.style={align=center}]
    \node[world] (w1)
          [label={left: \mFalse{p}},
            label={right:$[w_1]=[w_3]$},
            label={below: $\mSat{{}}{\Box p}$}] {$[w_1]$}; 
    \node[world] (w2)
         [label={right: \mTrue{p}},
           label = {below: $\mSat{{}}{\Box p}$},
           right=of w1, yshift=1.5cm] {$[w_2]$}; 
    \draw[->] (w1) to (w2);
    \node[world] (w4)
         [label={right:\mTrue{p}},
           label={below: $\mSat/{{}}{\Box p}$},
           right=of w1,yshift=-1.5cm] {$[w_4]$};
    \draw[->] (w1) to (w4);
    \draw[reflexive above] (w4) to (w4);     
    \node[world] (w5)
         [label={right: \mFalse{p}},
             label=below:$\mSat/{{}}{\Box p}$,
             right=of w4] {$[w_5]$}; 
    \draw[->, bend left] (w4) to (w5);
    \draw[->, bend left] (w5) to (w4);
    \draw[reflexive above] (w5) to (w5);     
  \end{tikzpicture}
  \caption{A filtração do modelo de \olref{fig:ser-eucl}.}
  \ollabel{fig:ser-eucl2}
\end{figure}

Em particular, para obter uma filtração euclidiana não basta considerar
filtrações por meio de conjuntos $\Gamma$ arbitrários fechados sob
sub-!!{formula}s. Precisamos considerar conjuntos $\Gamma$ que sejam
\emph{modalmente fechados} (veja \olref[pre]{defn:modallyclosed}). Tais
conjuntos de sentenças são infinitos e, portanto, não produzem de imediato
uma propriedade do modelo finito nem a decidibilidade do sistema
correspondente.

\begin{thm}\ollabel{thm:modal-closed-filt}
  Seja $\Gamma$ modalmente fechado, seja $\mModel{M}=\tuple{W,R,V}$ e
  seja $\mModel{M^*} = \tuple{W^*,R^*,V^*}$ uma filtração mais grossa de
  $\mModel{M}$.
  \begin{enumerate}
  \item Se $\mModel{M}$ é simétrico, $\mModel{M^*}$ também é.
  \item Se $\mModel{M}$ é transitivo, $\mModel{M^*}$ também é.
  \item Se $\mModel{M}$ é euclidiano, $\mModel{M^*}$ também é.
  \end{enumerate}
\end{thm}

\begin{proof}
\begin{enumerate}
  \item Exercício. Use o fato de que \Ax{B} e $\Ax{B_\Diamond}$ são
    válidas em todos os modelos simétricos.
  \item Se $\mModel{M^*}$ é uma filtração mais grossa, então, por
    definição, $R^*[u][v]$ vale se, e somente se, $C_1(u,v)$. Para a
    transitividade, suponha $C_1(u,v)$ e $C_1(v,w)$; precisamos mostrar
    $C_1(u,w)$. \iftag{prvBox}{Suponha $\mSat{M}{\Box !A}[u]$; então
      $\mSat{M}{\Box\Box!A}[u]$, pois \Ax{4} é válida em todos os
      modelos transitivos; como $\Box\Box!A \in \Gamma$ pelo fechamento,
      por $C_1(u,v)$ também $\mSat{M}{\Box!A}[v]$ e, por $C_1(v,w)$,
      também $\mSat{M}{!A}[w]$. }{}\iftag{prvDiamond}{Suponha
      $\mSat{M}{!A}[w]$; então $\mSat{M}{\Diamond !A}[v]$ por
      $C_1(v,w)$, pois $\Diamond !A \in \Gamma$ pelo fechamento modal.
      Por $C_1(u,v)$, obtemos $\mSat{M}{\Diamond\Diamond !A}[u]$, pois
      $\Diamond\Diamond!A \in \Gamma$ pelo fechamento modal. Como
      $\Ax{4}_\Diamond$ é válida em todos os modelos transitivos,
      $\mSat{M}{\Diamond!A}[u]$.}{}
  \item Exercício. Use o fato de que tanto \Ax{5} quanto $\Ax{5_\Diamond}$
    são válidas em todos os modelos euclidianos.
\end{enumerate}
\end{proof}

\begin{prob}
  Complete a demonstração de \olref[nml][fil][euc]{thm:modal-closed-filt}.
\end{prob}

\end{document}
